\section{Analysis}
\subsection{Ranking and accumulated weight diverge}
\begin{table}[t]
\centering\small\setlength{\tabcolsep}{3pt}
\begin{tabular}{lrrr}
\toprule
Pool / gallery / shot & Gallery coeff. (\%) & External share (\%) & Mean AUROC \\ \midrule
Original / Pets / 1 & 99.597 & 80.408 & 0.9280 \\
Original / DTD / 1 & 99.485 & 100.000 & -- \\
Original / Pets / 5 & 98.053 & 80.731 & 0.9477 \\
Original / DTD / 5 & 97.424 & 100.000 & -- \\
Confirm. / Pets / 1 & 99.605 & 80.826 & 0.9263 \\
Confirm. / DTD / 1 & 99.487 & 100.000 & -- \\
Confirm. / Pets / 5 & 98.071 & 80.862 & 0.9510 \\
Confirm. / DTD / 5 & 97.432 & 100.000 & -- \\
\bottomrule\end{tabular}
\caption{Final-update weight accounting. Gallery coefficient is the mean class-specific gallery fraction of support-plus-gallery mass. External share is the fraction of gallery weight on classes outside the episode. AUROC is undefined if no in-episode class is present. These descriptive coefficients are not causal directional effects.}\label{tab:influence}
\end{table}

Matched one-shot mean task AUROC is 0.9280 and 0.9263, but 80.408\% and 80.826\% of retained gallery weight is task-external. Ranking and retained mass answer different questions: the latter also depends on the number of irrelevant samples and their absolute weights. Good ordering does not demonstrate calibrated membership probabilities.

In mismatched one-shot conditions every gallery example is task-external, yet the final gallery coefficient averages approximately 99.5\%. Little scalar mass remains on support in the final unnormalized average. This is an algebraic observation, not a measure of feature-vector direction or a causal contribution for each image. All 4,000 task-condition predictions replay exactly, linking the accounting to the evaluated procedure.

\subsection{Mass and composition have distinct effects}
\begin{table}[t]
\centering\small\setlength{\tabcolsep}{3pt}
\begin{tabular}{lrrrrrr}
\toprule
Pool / gallery / shot & Original & Mass & Oracle & Oracle+mass & Zero & R2 \\ \midrule
Original / Pets / 1 & 94.92 & 95.76 & 97.92 & 96.94 & 94.75 & 96.24 \\
Original / DTD / 1 & 64.19 & 94.30 & 94.40 & 94.40 & 94.40 & 93.93 \\
Original / Pets / 5 & 97.19 & 98.59 & 98.79 & 98.83 & 98.44 & 98.44 \\
Original / DTD / 5 & 93.30 & 98.42 & 98.46 & 98.46 & 98.46 & 98.44 \\
Confirm. / Pets / 1 & 94.87 & 95.39 & 97.87 & 96.74 & 94.22 & 96.02 \\
Confirm. / DTD / 1 & 63.99 & 93.61 & 93.74 & 93.74 & 93.74 & 93.39 \\
Confirm. / Pets / 5 & 97.19 & 98.39 & 98.58 & 98.61 & 98.38 & 98.35 \\
Confirm. / DTD / 5 & 92.98 & 98.32 & 98.35 & 98.35 & 98.35 & 98.35 \\
\bottomrule\end{tabular}
\caption{Retrospective final-step interventions, accuracy in percent. Mass matches gallery weight to support mass per class. Oracle removes examples whose true class is outside the episode; Oracle+mass combines both. Oracle variants require privileged labels. Earlier fitted weights and centering stay fixed.}\label{tab:interventions}
\end{table}

Final mass balancing raises mismatched one-shot accuracy from 64.192\% to 94.301\% in the original pool and from 63.992\% to 93.613\% in confirmation. Most mismatch harm can be removed by this constrained intervention. However, matched one-shot accuracy remains below R2 by 0.483 and 0.624 pp, with both paired intervals negative (Appendix~\ref{app:interventions}). Mass correction alone does not meet the improvement objective.

Oracle filtering gives matched one-shot accuracies of 97.923\% and 97.867\%, exceeding R2 by 1.683 and 1.851 pp. This diagnoses composition in an already-fitted state; it is not an admissible deployment comparison because gallery labels are unavailable at inference. It is also not a theoretical upper bound: earlier contaminated updates remain fixed. Combining filtering and balancing changes performance again, so the two operations are not interchangeable.

Both diagnostics reuse the main evaluation tasks rather than adding independent test samples. Weight accounting is descriptive; final-step interventions isolate particular changes to a saved fitted state. Neither establishes that a learned membership estimator will generalize or that membership is the only useful next design direction.
